\begin{theorem}[有限模约化阴影的完全边际分类]\label{thm:xi-crt-invisibility-marginal-prime-spectrum}
设
\[
A=\bigoplus_{i=1}^{r}\mathbb Z[S_i^{-1}],
\]
其中每个 \(S_i\) 都是有限素数集。
对每个素数 \(p\) 定义边际局部化计数
\[
m_p:=\card{\{i\in\{1,\dots,r\}:p\in S_i\}}.
\]
对任意整数 \(n\ge 1\) 写其素因子分解 \(n=\prod_{p}p^{a_p}\)。
则有自然同构
\[
\boxed{\;
A/nA\ \cong\ \prod_{p\mid n}\Bigl(\mathbb Z/p^{a_p}\mathbb Z\Bigr)^{\,r-m_p}.\;}
\]
并且整个有限模约化族 \(\{A/nA\}_{n\ge 1}\) 恰好决定且仅决定数据
\[
r,\qquad (m_p)_{p\in\mathbb P}.
\]
更具体地，
\[
\boxed{\;
r=\sup_{p\ \mathrm{prime}}\dim_{\mathbb F_p}(A/pA),
\qquad
m_p=r-\dim_{\mathbb F_p}(A/pA).\;}
\]
因此任何更高阶的共局部化相关结构，都不可能由全部有限模约化阴影恢复。
\leanverified{paper\_xi\_crt\_invisibility\_marginal\_prime\_spectrum}
\end{theorem}
\begin{proof}
固定 \(i\)。
令
\[
n_i:=\prod_{p\mid n,\ p\notin S_i}p^{a_p}.
\]
当 \(p\in S_i\) 时，\(p\) 在 \(\mathbb Z[S_i^{-1}]\) 中可逆，故 \(p^{a_p}\mathbb Z[S_i^{-1}]=\mathbb Z[S_i^{-1}]\)，从而
\[
n\,\mathbb Z[S_i^{-1}]=n_i\,\mathbb Z[S_i^{-1}].
\]
又因 \(n_i\) 的所有素因子都不属于 \(S_i\)，局部化不会改变模 \(n_i\) 的商，故
\[
\mathbb Z[S_i^{-1}]/n\,\mathbb Z[S_i^{-1}]
\cong
\mathbb Z/n_i\mathbb Z.
\]
于是
\[
A/nA
\cong
\prod_{i=1}^{r}\mathbb Z/n_i\mathbb Z.
\]
对每个 \(i\)，由中国剩余定理，
\[
\mathbb Z/n_i\mathbb Z
\cong
\prod_{p\mid n,\ p\notin S_i}\mathbb Z/p^{a_p}\mathbb Z.
\]
将所有分量按素数 \(p\) 重组，即得
\[
A/nA
\cong
\prod_{p\mid n}\Bigl(\mathbb Z/p^{a_p}\mathbb Z\Bigr)^{\card\{i:\ p\notin S_i\}}
=
\prod_{p\mid n}\Bigl(\mathbb Z/p^{a_p}\mathbb Z\Bigr)^{\,r-m_p}.
\]

接着说明反向恢复。
当 \(n=p\) 为素数时，上式退化为
\[
A/pA\cong (\mathbb Z/p\mathbb Z)^{\,r-m_p},
\]
故
\[
\dim_{\mathbb F_p}(A/pA)=r-m_p.
\]
由于每个 \(S_i\) 仅含有限多个素数，故只有有限多个 \(p\) 满足 \(m_p\neq 0\)；
于是对充分大的素数 \(p\) 有 \(m_p=0\)，从而
\[
\dim_{\mathbb F_p}(A/pA)=r.
\]
因此
\[
r=\sup_{p\ \mathrm{prime}}\dim_{\mathbb F_p}(A/pA),
\qquad
m_p=r-\dim_{\mathbb F_p}(A/pA).
\]
这就说明 \(\{A/nA\}_{n\ge 1}\) 唯一恢复 \(r\) 与全部 \(m_p\)。
证毕。
\end{proof}

\begin{corollary}[高阶共局部化张量的有限模不可见性]\label{cor:xi-crt-indistinguishable-pair}
\leanverified{paper\_xi\_crt\_indistinguishable\_pair}
对任意有限素数集
\[
T=\{p_1,\dots,p_m\},
\qquad
m\ge 2,
\]
定义高阶共局部化计数
\[
J_T(A):=\card{\{i:\ T\subseteq S_i\}}.
\]
则 \(J_T\) 不是有限模约化族 \(\{A/nA\}_{n\ge 1}\) 的函数。
更具体地，令
\[
A_T:=\bigoplus_{j=1}^{m}\mathbb Z[1/p_j],
\qquad
A_T':=\mathbb Z[(p_1\cdots p_m)^{-1}]\oplus \mathbb Z^{m-1}.
\]
则对一切 \(n\ge 1\) 都有
\[
\boxed{\;A_T/nA_T\cong A_T'/nA_T'.\;}
\]
但
\[
\boxed{\;J_T(A_T)=0,\qquad J_T(A_T')=1.\;}
\]
因此有限模约化阴影严格抹去全部 \(|T|\ge 2\) 的共现几何，只保留单点边际数据 \((m_p)_p\)。
\end{corollary}
\begin{proof}
对 \(A_T\)，每个素数 \(p_j\) 恰出现于一个分量 \(\mathbb Z[1/p_j]\) 中，故
\[
m_{p_j}(A_T)=1
\qquad(1\le j\le m),
\]
其余素数的边际计数为零，总秩为 \(m\)。

对 \(A_T'\)，唯一局部化分量 \(\mathbb Z[(p_1\cdots p_m)^{-1}]\) 同时包含全部 \(p_j\)，其余 \(m-1\) 个分量为纯整数分量。
因此同样有
\[
m_{p_j}(A_T')=1
\qquad(1\le j\le m),
\]
其余边际计数仍为零，总秩仍为 \(m\)。
由定理 \ref{thm:xi-crt-invisibility-marginal-prime-spectrum}，遂得
\[
A_T/nA_T\cong A_T'/nA_T'
\qquad(\forall n\ge 1).
\]

另一方面，在 \(A_T\) 中没有任何单一分量同时包含全部 \(p_1,\dots,p_m\)，故 \(J_T(A_T)=0\)；
而在 \(A_T'\) 中，第一个分量恰同时包含全部这些素数，故 \(J_T(A_T')=1\)。
证毕。
\end{proof}

\endinput
